% Copyright (c) 2026 Eugene Joseph M. Ragasa.
% SPDX-License-Identifier: Apache-2.0
% Author: Eugene Joseph M. Ragasa
% Formal proof completed: 2026-10-31
\documentclass[10pt]{article}
\usepackage[T1]{fontenc}
\usepackage{amsmath,amssymb,booktabs,tabularx}
\usepackage[margin=0.75in]{geometry}
\usepackage[hidelinks]{hyperref}
\newcolumntype{Y}{>{\raggedright\arraybackslash}X}
\title{Proof Dossier: Equal-Phase Power Law on a Normalized Product Torus}
\author{Eugene Joseph M. Ragasa}
\date{Formal proof completed: 2026-10-31}
\begin{document}
\raggedright
\maketitle

\section{Identity and status}
\begin{tabularx}{\textwidth}{@{}lY@{}}
\toprule
Statement ID & \texttt{BG-FOURIER-THM-125} \\
Status & proved; no admissions \\
Lean declaration & \path{ProjectKoios.BandGauge.Fourier.normalizedProductTorusAverage_exp_samePhase_eq_pow} \\
Owning module & \path{ProjectKoios/BandGauge/Fourier/NormalizedProductTorus.lean} \\
Author and license & Eugene Joseph M. Ragasa; \texttt{Apache-2.0} \\
\bottomrule
\end{tabularx}

\section{Exact mathematical statement}
Let $D$ be finite, let $\mathbb T_1=\mathbb R/\mathbb Z$ carry its standard Haar probability
measure $\mu$, and let $\varphi:\mathbb T_1\to\mathbb C$ be one supplied phase. Then
\[
 \int_{\mathbb T_1^D}
   \exp\!\left(\sum_{d\in D}\varphi(q_d)\right)\,d\mu^D(q)
 =\left(\int_{\mathbb T_1}\exp(\varphi(t))\,d\mu(t)\right)^{|D|}.
\]
For $D=\operatorname{Fin}(2)$ and $D=\operatorname{Fin}(3)$, separately named checked
corollaries give the square and cube.

\section{Source relation and immutable identity}
The orthogonal KSD 2-D flat-parent protocol states the squared comparator at lines 414--435, and
the orthogonal 3-D protocol states the factorized cubed comparator at lines 359--381. The source
identity is KSD commit \texttt{fcfc4d5fa76b08c040a02c97a109eb5f5fc06e48}, tree
\texttt{fe8b47b33fcf0b24c5b28f8799ed220e3f6d7d20}. The exact source paths, SHA-256 fingerprints,
read-only access mode, and dirty-state limitation are recorded in \path{docs/source-audit.md}.

The Formal base is PR~\#44 merge commit
\texttt{24263f08a5b69fc34be3073bc724a3b0ad4b80d6}, tree
\texttt{2cbef2a05174f5ad491392b40dd4265a751ed932}. The owning-module SHA-256 is
\texttt{a9d7a1fb58123cbeeb8e8474b6e8b54d847d59a962945fe09db3f84a42c746c8}.
No topic commit is claimed before checkpoint authorization.

\section{Assumptions and conventions}
The theorem uses a finite coordinate type, Mathlib's unit additive circle in every coordinate, the
finite product of its standard total-mass-one Haar measure, and one exact complex phase reused in
every coordinate. The dimension is the exact cardinality of $D$. The theorem does not infer that
an external phase family has equal coordinate factors.

\section{Prior art and contribution boundary}
The result is a standard equal-factor specialization of finite-product Fubini, not a novelty claim.
The Lean proof specializes pinned Mathlib theorem
\texttt{MeasureTheory.integral\_fintype\_prod\_eq\_pow}. The local contribution is a typed,
normalized BandGauge statement with explicit dimension-cardinality, square, cube, and evidence
boundaries suitable for later source-specific bridges.

\section{Proof structure}
Use $\exp(\sum_d\varphi(q_d))=\prod_d\exp(\varphi(q_d))$. Finite-product Fubini turns the
joint average into the product of one-circle averages. Since every coordinate factor is identical,
the finite product is the corresponding scalar power with exponent $|D|$.

\section{Lean realization and trusted base}
The proof uses \texttt{Complex.exp\_sum} and the equal-factor power theorem in
\path{Mathlib/MeasureTheory/Integral/Pi.lean}, lines 119--126. Toolchain: Lean
\texttt{4.35.0} at \texttt{eb6cafb7290e8bd5f2937522d8d804f58f02bafc}; Mathlib
\texttt{56a005855ea17e0b0958f73ec801ae140b0a1679}.
\begin{verbatim}
#print axioms ProjectKoios.BandGauge.Fourier.
  normalizedProductTorusAverage_exp_samePhase_eq_pow
-- [propext, Classical.choice, Quot.sound]
\end{verbatim}
No repository-owned axiom or admission occurs.

\section{Evidence and interpretation boundaries}
The right-hand one-circle integral is deliberately unevaluated. The theorem does not identify it
with $J_0(a)$ or $J_0(2a/3)$, prove the KSD phase or scaling, certify a continuum comparator,
authenticate a quadrature or mesh, or bound numerical error. The exact square and cube corollaries
do not transfer scientific, certificate, or publication status.

\section{AI, automation, and review record}
The Lean proof and dossier were AI-assisted and checked by the maintainer. Kernel checking is not
independent mathematical, scientific, source-fidelity, novelty, or publication review.

\section{Unresolved consumer obligations}
A consumer must prove the coordinate phase identification and one-circle integral under its frozen
coordinate convention, then separately verify any numerical approximation, source artifact,
tolerance, and retained result. A Bessel theorem, if later adopted, requires its own source and
statement audit.

\section{Review and revision history}
\begin{tabularx}{\textwidth}{@{}lllY@{}}
\toprule Date & Revision & Process & Outcome and limitation \\
\midrule 2026-10-31 & initial dossier & AI-assisted maintainer review & exact cardinality-power law proved; one-circle evaluation and external binding remain open \\
\bottomrule
\end{tabularx}
\end{document}
